\documentclass[11pt,letterpaper]{article}
\usepackage[T1]{fontenc}
\usepackage[utf8]{inputenc}
\usepackage{lmodern,microtype}
\usepackage[margin=0.79in,headheight=15pt]{geometry}
\usepackage{parskip,enumitem,longtable,tabularx,booktabs,array,xcolor,fancyhdr,hyperref}
\definecolor{navy}{HTML}{17314C}
\definecolor{teal}{HTML}{087F78}
\hypersetup{colorlinks=true,linkcolor=navy,urlcolor=teal}
\setlength{\emergencystretch}{3em}
\setlist[itemize]{leftmargin=*,nosep,topsep=5pt}
\setlist[enumerate]{leftmargin=*,itemsep=3pt}
\renewcommand{\arraystretch}{1.22}
\pagestyle{fancy}\fancyhf{}
\fancyhead[L]{\small Enterprise Data Modeling}\fancyhead[R]{\small Documentation Suite}
\fancyfoot[C]{\thepage}
\newcommand{\srcbook}{\textit{The Data Model Resource Book}, revised edition, Volume 2}
\newcommand{\ResponseLine}{\par\vspace{1.05em}\noindent\rule{\linewidth}{0.4pt}\par}
\newcommand{\ResponseArea}{\ResponseLine\ResponseLine\ResponseLine}
\begin{document}
\begin{center}
{\LARGE\bfseries\color{navy} Normalization and Dependency Analysis}\par
\vspace{0.25em}{\large Guided workbook | conceptual and operational analysis}\par
\vspace{0.3em}{\small Original engineering study material \textbullet{} 2026 edition}
\end{center}
\vspace{0.8em}
\section{Workbook objective}
Complete independently before looking at any model answer. Record assumptions, counterexamples and a proof of each chosen invariant. This is supplementary material; no exercise copies a worked example from the source book.


\section{Problem A: supplier orders}
Given: R(OrderID, LineNo, SupplierID, SupplierName, ProductID, ProductName, Quantity). Candidate key (OrderID, LineNo). Assume OrderID determines SupplierID; SupplierID determines SupplierName; ProductID determines ProductName; (OrderID, LineNo) determines ProductID and Quantity.

Tasks: Find all nontrivial functional dependencies you need; state 1NF assumptions; exhibit partial and transitive dependencies; decompose to 3NF and identify a join that demonstrates losslessness.

\subsection{Your working}
\ResponseArea\ResponseArea
\subsection{Verification notes}
\ResponseArea

\section{Problem B: overlapping identities}
Given: R(PartyID, RoleCode, EffectiveFrom, EffectiveTo, Name). A party has multiple roles and may change its name. The business prohibits overlapping intervals for the same role and party.

Tasks: Define the candidate key and interval convention; explain why standard row-local CHECK constraints cannot detect arbitrary conflicting rows; design a concurrency-safe exclusion strategy or transactional validation.

\subsection{Your working}
\ResponseArea\ResponseArea
\subsection{Verification notes}
\ResponseArea

\section{Problem C: workload trade-off}
Given: A reporting query repeatedly aggregates order line totals by customer. The normalized schema is correct but performance is poor.

Tasks: Propose index, materialized aggregation and precomputation alternatives; specify freshness requirements, recomputation and a reconciliation test.

\subsection{Your working}
\ResponseArea\ResponseArea
\subsection{Verification notes}
\ResponseArea

\section{Self-assessment}
\begin{itemize}
\item Can another reviewer reproduce the decomposition or integrity argument?
\item Did the model document what would falsify its assumptions?
\item Can every proposed invariant be linked to executable or review evidence?
\end{itemize}

\vfill
\hrule
{\footnotesize\textbf{Provenance.} Supplementary engineering practice: original content; the cited book provides domain-modeling context.}
\end{document}
